\section{Временная диаграмма выполнения заданий}

На рисунке~\ref{fig:timeline} представлена временная диаграмма выполнения практических работ. Синими точками отмечены дни, в которые выполнялись коммиты в репозиторий; число над точкой показывает количество коммитов за день. Красными точками отмечены даты сдачи или состояние защиты практических работ.

\begin{figure}[H]
    \centering
    \begin{tikzpicture}[x=0.78cm, y=1cm, font=\scriptsize]
        \draw[->, thick] (0,0) -- (15.2,0);

        \foreach \x/\date/\count in {
            0/27.02/1,
            1/05.03/1,
            2/06.03/4,
            3/13.03/3,
            4/19.03/1,
            5/20.03/5,
            6/27.03/2,
            7/20.04/8,
            8/23.04/3,
            9/24.04/6,
            10/08.05/4,
            11/15.05/13,
            12/21.05/5,
            13/22.05/14,
            14/23.05/1
        } {
            \fill[blue] (\x,0.45) circle (5pt);
            \node[above] at (\x,0.58) {\count};
            \node[below] at (\x,-0.08) {\date};
        }

        \fill[red] (5,-0.62) circle (5pt);
        \node[below, align=center] at (5,-0.78) {20.03\\лаб.1 зачтена};

        \fill[red] (9,-0.62) circle (5pt);
        \node[below, align=center] at (9,-0.78) {24.04\\лаб.2 зачтена};

        \fill[red] (13,-0.62) circle (5pt);
        \node[below, align=center] at (13,-0.78) {22.05\\лаб.3 зачтена\\лаб.4 один вопрос};

        \fill[blue] (0,-2.05) circle (6pt);
        \node[right, align=left] at (0.25,-2.05) {Коммиты: один кружок на день,\\указано количество коммитов};
        \fill[red] (0,-2.75) circle (6pt);
        \node[right, align=left] at (0.25,-2.75) {Сдачи практических работ: дата\\и результат защиты};
    \end{tikzpicture}
    \caption{Временная диаграмма коммитов и сдач практических работ}
    \label{fig:timeline}
\end{figure}

\newpage

\section{Структура и содержание разделов репозитория}

Репозиторий с практическими работами имеет следующую структуру:

\begin{verbatim}
practice/
|-- practice_01/
|-- practice_02/
|   |-- MyEvolProject/
|   \-- myParser/
|-- practice_03/
|   |-- glitcher/
|   \-- maze-rws/
|-- practice_04/
|   \-- mazeGame/
|       |-- app/
|       |-- data/
|       \-- src/
|-- practice_05/
|   \-- qc/
|       |-- app/
|       |-- src/
|       \-- test/
\-- report/
    |-- main.tex
    |-- main.pdf
    \-- tex/
\end{verbatim}

Каталог \mintinline{text}{practice_01} содержит первые функции и решения на защиту. В \mintinline{text}{practice_02} находятся проекты с пользовательскими типами и парсерами. В \mintinline{text}{practice_03} расположены проекты \mintinline{text}{glitcher} и \mintinline{text}{maze-rws}. В \mintinline{text}{practice_04/mazeGame} находится текстовая игра по лабиринту с собственным трансформером состояния. В \mintinline{text}{practice_05/qc} находится проект с QuickCheck-тестами. Все изменения фиксировались в системе контроля версий Git.

\newpage
